% संपूर्ण मराठी ग्रंथातील प्रकरण 54: निर्दोषता आणि संपूर्णता
% हा संपादनयोग्य स्रोतखंड आहे. संकलनासाठी संपूर्ण स्रोतसंग्रहातील
% अक्षररूपे, पूर्वभाग, चिन्हव्याख्या आणि आकृती-स्रोत आवश्यक आहेत.
% मूळ: Open Logic Project; मजकूर आणि रूपांतर CC BY 4.0.
% यंत्रानुवाद, दुरुस्ती आणि तपासणी: OpenAI Codex —
% GPT-5.6 Sol आणि GPT-6 Sol, दोन्ही Ultra effort.
% दुरुस्ती-पुनरावलोकन: GPT-6.1 Sol, Ultra effort.
\chapter{निर्दोषता आणि संपूर्णता}\label{int:sc:chap}
\begin{editorial}
  या प्रकरणात अंतःप्रज्ञावादी विधानीय तर्कशास्त्रासाठी
  निर्दोषता आणि संपूर्णता यांविषयीचे निष्कर्ष
  एकत्र दिले आहेत. याला प्रस्तावना हवी आहे.
  संपूर्णतेची सिद्धता सिद्धतायोग्यतेबद्दलच्या
  अशा तथ्यांचा आधार घेते की जी कुठेतरी
  स्पष्टपणे मांडून सिद्ध करायला हवीत.
\end{editorial}









\section{स्वयंसिद्धकीय निष्पत्तीची निर्दोषता}\label{int:sc:sax:sec}

\begin{editorial}
  निर्दोषतेची सिद्धता सर्व स्वयंसिद्धके
  अंतःप्रज्ञावादी अर्थाने वैध आहेत, या
  तथ्यावर अवलंबून आहे. हे अजून सिद्ध
  करायचे आहे; उदाहरणार्थ, चिन्हार्थमीमांसेच्या
  प्रकरणात ते करता येईल.
\end{editorial}

\begin{thm}[निर्दोषता]
  \label{int:sc:sax:thm:soundness}
  $\Gamma \Proves \varphi$ असेल, तर
  $ \Gamma \Entails \varphi$.
\end{thm}

\begin{proof}
  $\Gamma \Proves \varphi$ असेल, तर
  $\Gamma \Entails \varphi$, हे सिद्ध करू.
  $\Gamma$ पासून~$\varphi$ च्या
  निष्पत्तीमधील सूत्रांची
  संख्या~$n$ यावर विगमनाने सिद्धता करतो.
  $\varphi_1$, \dots, $\varphi_n = \varphi$ ही
  $\Gamma$ पासूनची निष्पत्ती असेल,
  तर $\Gamma \Entails \varphi_n$ हे दाखवू.
  $\varphi_1$, \dots, $\varphi_n$ ही
  निष्पत्ती असेल, तर कोणत्याही $k<n$
  साठी $\varphi_1$, \dots, $\varphi_k$ हीही
  निष्पत्ती आहे, हे लक्षात घ्या.

  लांबी~$0$ च्या निष्पत्त्या नसतात; म्हणून
  $n=0$ साठी दावा रिक्तपणे खरा आहे.
  त्यामुळे लांबी~$<n$ असलेल्या सर्व
  निष्पत्त्या साठी दावा खरा आहे, असे
  विगमन-गृहीतक घेऊ. आता~$\varphi_n$ चे
  समर्थन कशावर आहे त्यानुसार प्रकार पाहू.

  \begin{enumerate}
  \item $\varphi_n$ हे स्वयंसिद्धक आहे.
    सर्व स्वयंसिद्धके वैध असल्याने कोणत्याही~$\Gamma$
    साठी $\Gamma \Entails
    \varphi_n$.
  \item $\varphi_n \in \Gamma$. मग कोणतेही
    $\mModel{M}$ आणि $w$ घ्या.
    $\mSat{M}{\Gamma}[w]$ असेल, तर उघडच
    $\mSat{M}{\varphi_n}[w]$; म्हणजे
    $\Gamma \Entails \varphi$.
  \item $\varphi_n$ हे $\varphi_i$ आणि
    $\varphi_j \ident \varphi_i \lif
    \varphi_n$ यांपासून \MP{} ने निष्पन्न होते.
    $\varphi_1$, \dots, $\varphi_i$ आणि
    $\varphi_1$, \dots, $\varphi_j$ या
    $\Gamma$ पासूनच्या निष्पत्त्या आहेत.
    म्हणून विगमन-गृहीतकानुसार
    $\Gamma \Entails \varphi_i$ आणि
    $\Gamma \Entails \varphi_i \lif \varphi_n$.

    $\mSat{M}{\Gamma}[w]$ असे समजा.
    $\mSat{M}{\Gamma}[w]$ आणि
    $\Gamma \Entails \varphi_i \lif \varphi_n$
    असल्यामुळे $\mSat{M}{\varphi_i \lif
      \varphi_n}[w]$. व्याख्येनुसार याचा अर्थ
    असा की $Rww'$ असलेल्या प्रत्येक $w'$
    साठी $\mSat{M}{\varphi_i}[w']$ असेल, तर
    $\mSat{M}{\varphi_n}[w']$. $R$ परावर्ती
    असल्याने $w$ हेही $Rww'$ पूर्ण करणाऱ्या
    $w'$ पैकी एक आहे. म्हणजे
    $\mSat{M}{\varphi_i}[w]$ असेल, तर
    $\mSat{M}{\varphi_n}[w]$.
    $\Gamma \Entails \varphi_i$ असल्याने
    $\mSat{M}{\varphi_i}[w]$. म्हणून
    अपेक्षेप्रमाणे $\mSat{M}{\varphi_n}[w]$.
  \end{enumerate}
\end{proof}












\section{नैसर्गिक निगमनाची निर्दोषता}\label{int:sc:snd:sec}

आता संबंधी चिन्हार्थमीमांसेच्या दृष्टीने नैसर्गिक निगमनाची
निर्दोषता सिद्ध करू. म्हणजेच, गृहीतकांच्या संचापासून
एखादे सूत्र निष्पन्न करता येण्याजोगे असेल, तर त्या
गृहीतकांच्या संचापासून ते सूत्र
तार्किकरीत्या निष्पन्न होते, हे दाखवू.

\begin{thm}[निर्दोषता]
  \label{int:sc:snd:thm:soundness}
  $\Gamma \Proves \varphi$ असेल, तर $ \Gamma \Entails \varphi$.
\end{thm}

\begin{proof}
  $\Gamma \Proves \varphi$ असेल, तर $\Gamma \Entails \varphi$,
  हे सिद्ध करू. $\Gamma$ पासून~$\varphi$ च्या
  निष्पत्ती वर विगमनाने सिद्धता करतो.

  \begin{enumerate}
  \item निष्पत्तीमध्ये केवळ गृहीतक~$\varphi$ असेल,
    तर $\varphi \Proves \varphi$. आता $\varphi \Entails \varphi$ हे
    दाखवायचे आहे. $ \mSat{M}{\varphi}[w]$ असे समजा.
    मग उघडच $\mSat{M}{\varphi}[w]$.

  \item निष्पत्तीचा शेवट $\Intro{\land}$ ने होतो:
    आधारविधान~$\psi$ ची
      मुक्त न केलेले गृहीतके~$\Gamma$ यांपासून आणि
      आधारविधान~$\chi$ ची मुक्त न केलेले
      गृहीतके~$\Delta$ यांपासून असलेल्या
      निष्पत्त्यांवरून $\Gamma \Proves
      \psi$ आणि $\Delta \Proves \chi$ हे मिळते.
      विगमन-गृहीतकानुसार $\Gamma \Entails \psi$
      आणि $\Delta \Entails \chi$. संपूर्ण
      निष्पत्तीची मुक्त न केलेले गृहीतके
      $\Gamma$ व $\Delta$ मिळून असल्यामुळे
      $\Gamma \cup \Delta \Entails \psi \land \chi$
      हे दाखवायचे आहे. म्हणून
      $\mSat{M}{\Gamma \cup \Delta}[w]$ असे समजा.
      मग $\mSat{M}{\Gamma}[w]$ सुद्धा.
      $\Gamma \Entails \psi$ असल्याने
      $\mSat{M}{\psi}[w]$. त्याचप्रमाणे
      $\mSat{M}{\chi}[w]$. म्हणून
      $\mSat{M}{\psi \land \chi}[w]$.

  \item निष्पत्तीचा शेवट $\Elim{\land}$ ने होतो:
    आधारविधान $\psi
        \land \chi$ ची मुक्त न केलेले गृहीतके
        $\Gamma$ यांपासून असलेली निष्पत्ती
        $\Gamma \Proves \psi \land \chi$ हे दाखवते.
        विगमन-गृहीतकानुसार
        $\Gamma \Entails \psi \land \chi$.
        $\Gamma \Entails \psi$ हे दाखवायचे आहे.
        म्हणून $\mSat{M}{\Gamma}[w]$ असे समजा.
        $\Gamma \Entails \psi \land \chi$
        असल्याने $\mSat{M}{\psi \land \chi}[w]$.
        मग $\mSat{M}{\psi}[w]$ सुद्धा.
        त्याचप्रमाणे $\Elim\land$ चा शेवट~$\chi$
        ने होत असेल, तर $\Gamma \Entails \chi$.

  \item निष्पत्तीचा शेवट $\Intro{\lor}$ ने होतो:
    आधारविधान $\psi$ आहे,
      आणि~$\psi$ येथे संपणाऱ्या निष्पत्तीची
      मुक्त न केलेले गृहीतके $\Gamma$ आहेत,
      असे समजा. मग $\Gamma \Proves \psi$;
      विगमन-गृहीतकानुसार $\Gamma \Entails \psi$.
      $\Gamma \Entails \psi \lor \chi$ हे दाखवायचे आहे.
      $\mSat{M}{\Gamma}[w]$ असे समजा.
      $\Gamma \Entails \psi$ असल्याने
      $\mSat{M}{\psi}[w]$. पण मग
      $\mSat{M}{\psi \lor \chi}[w]$ सुद्धा.
      त्याचप्रमाणे आधारविधान~$\chi$ असेल, तर
      $\Gamma \Entails \chi$.

  \item निष्पत्तीचा शेवट~$\Elim{\lor}$ ने होतो:
    आधारविधानांवर संपणाऱ्या
      निष्पत्त्या अनुक्रमे $\psi \lor \chi$ ची
      मुक्त न केलेले गृहीतके~$\Gamma$,
      $\theta$ ची मुक्त न केलेले गृहीतके
      $\Delta_1 \cup \{\psi\}$, आणि $\theta$ ची
      मुक्त न केलेले गृहीतके
      $\Delta_2 \cup \{\chi\}$ यांपासून आहेत.
      म्हणून $\Gamma \Proves \psi \lor \chi$,
      $\Delta_1 \cup \{\psi\} \Proves \theta$,
      आणि $\Delta_2 \cup \{\chi\} \Proves \theta$.
      विगमन-गृहीतकानुसार
      $\Gamma \Entails \psi \lor \chi$,
      $\Delta_1 \cup \{\psi\} \Entails \theta$,
      आणि $\Delta_2 \cup \{\chi\} \Entails \theta$.
      $\Gamma \cup \Delta_1 \cup \Delta_2 \Entails \theta$
      हे सिद्ध करायचे आहे.

    $\mSat{M}{\Gamma \cup \Delta_1 \cup \Delta_2}[w]$
    असे समजा. मग $\mSat{M}{\Gamma}[w]$;
    $\Gamma \Entails \psi \lor \chi$ असल्याने
    $\mSat{M}{\psi \lor \chi}[w]$.
    $\mSat{M}{}$ च्या व्याख्येनुसार
    $\mSat{M}{\psi}[w]$ किंवा
    $\mSat{M}{\chi}[w]$. म्हणून दोन प्रकार पाहू:
    (a)~$\mSat{M}{\psi}[w]$. मग
    $\mSat{M}{\Delta_1 \cup \{\psi\}}[w]$.
    $\Delta_1 \cup \{\psi\} \Entails \theta$
    असल्याने $\mSat{M}{\theta}[w]$.
    (b)~$\mSat{M}{\chi}[w]$. मग
    $\mSat{M}{\Delta_2 \cup \{\chi\}}[w]$.
    $\Delta_2 \cup \{\chi\} \Entails \theta$
    असल्याने $\mSat{M}{\theta}[w]$.
    म्हणून दोन्ही प्रकारांत
    $\mSat{M}{\theta}[w]$, हेच दाखवायचे होते.

  \item निष्पत्तीचा शेवट $\Intro{\lif}$ ने होतो
    आणि निष्कर्ष~$\psi\lif \chi$ आहे. मग
    आधारविधान~$\chi$ आहे, आणि त्या आधारविधानावर
    संपणाऱ्या निष्पत्तीची
    मुक्त न केलेले गृहीतके~$\Gamma \cup
    \{\psi\}$ आहेत. म्हणून
    $\Gamma \cup \{\psi\} \Proves \chi$;
    विगमन-गृहीतकानुसार
    $\Gamma \cup \{\psi\} \Entails \chi$.
    $\Gamma \Entails \psi \lif \chi$ हे दाखवायचे आहे.

    $\mSat{M}{\Gamma}[w]$ असे समजा.
    $Rww'$ असे असलेल्या सर्व $w'$ साठी
    $\mSat{M}{\psi}[w']$ असेल, तर
    $\mSat{M}{\chi}[w']$, हे दाखवायचे आहे.
    म्हणून $Rww'$ आणि $\mSat{M}{\psi}[w']$
    असे गृहीत धरा.
    \ref{int:sem:rel:prop:true-monotonic} नुसार
    $\mSat{M}{\Gamma}[w']$.
    $\Gamma \cup \{\psi\} \Entails \chi$
    असल्याने $\mSat{M}{\chi}[w']$,
    हेच दाखवायचे होते.

  \item निष्पत्तीचा शेवट $\Elim{\lif}$ ने होतो
    आणि निष्कर्ष~$\chi$ आहे. आधारविधाने
    $\psi \lif \chi$ आणि $\psi$ आहेत;
    त्यांची निष्पत्त्या अनुक्रमे
    मुक्त न केलेले गृहीतके $\Gamma$ आणि
    $\Delta$ यांपासून आहेत. म्हणून
    $\Gamma \Proves \psi \lif \chi$ आणि
    $\Delta \Proves \psi$.
    विगमन-गृहीतकानुसार
    $\Gamma \Entails \psi \lif \chi$
    आणि $\Delta \Entails \psi$.
    $\Gamma \cup \Delta \Entails \chi$
    हे दाखवायचे आहे.

    $\mSat{M}{\Gamma \cup \Delta}[w]$ असे समजा.
    $\mSat{M}{\Gamma}[w]$ आणि
    $\Gamma \Entails \psi \lif \chi$ असल्याने
    $\mSat{M}{\psi \lif \chi}[w]$.
    व्याख्येनुसार, $Rww'$ असे असलेल्या
    सर्व $w'$ साठी $\mSat{M}{\psi}[w']$
    असेल, तर $\mSat{M}{\chi}[w']$.
    $R$ परावर्ती असल्याने $w$ हे
    $Rww'$ असलेल्या $w'$ पैकी एक आहे.
    म्हणजे $\mSat{M}{\psi}[w]$ असेल,
    तर $\mSat{M}{\chi}[w]$.
    $\mSat{M}{\Delta}[w]$ आणि
    $\Delta \Entails \psi$ असल्याने
    $\mSat{M}{\psi}[w]$. म्हणून
    अपेक्षेप्रमाणे $\mSat{M}{\chi}[w]$.

  \item निष्पत्तीचा शेवट $\FalseInt$ ने होतो
    आणि निष्कर्ष~$\varphi$ आहे. आधारविधान
    $\lfalse$ आहे; त्या आधारविधानाची
    निष्पत्ती ज्या
    मुक्त न केलेले गृहीतकांपासून आहे,
    ती~$\Gamma$ आहेत. मग
    $\Gamma \Proves \lfalse$.
    विगमन-गृहीतकानुसार
    $\Gamma \Entails \lfalse$.
    $\Gamma \Entails \varphi$ हे दाखवायचे आहे.

    अप्रत्यक्ष सिद्धता करू.
    $\Gamma \Entails/ \varphi$ असेल, तर
    असे प्रतिमान~$\mModel{M}$ आणि जग~$w$
    आहेत की $\mSat{M}{\Gamma}[w]$
    आणि $\mSat/{M}{\varphi}[w]$.
    $\Gamma \Entails \lfalse$
    असल्याने $\mSat{M}{\lfalse}[w]$.
    पण हे अशक्य आहे; कारण व्याख्येनुसार
    $\mSat/{M}{\lfalse}[w]$.
    म्हणून $\Gamma \Entails \varphi$.
  \item निष्पत्तीचा शेवट $\Intro\lnot$ ने होतो: सराव.
  \item निष्पत्तीचा शेवट $\Elim\lnot$ ने होतो: सराव.
  \end{enumerate}
\end{proof}

\begin{prob}
  \ref{int:sc:snd:thm:soundness} ची सिद्धता
  पूर्ण करा. $\Intro\lnot$ आणि $\Elim\lnot$
  या प्रकारांसाठी
  \ref{int:sem:rel:defn:true-at-w} मधील
  $\mSat{M}{\lnot \varphi}[w]$ ची व्याख्या वापरा;
  म्हणजे $\lnot \varphi$ हे
  $\varphi \lif \lfalse$ ने व्याख्यित आहे
  असे मानू नका.
\end{prob}

\begin{prob}
  खालील सूत्रे अंतःप्रज्ञावादी तर्कशास्त्रात
  निष्पन्न करता येण्याजोगे नाहीत, हे दाखवा:
  \begin{enumerate}
    \item $(\varphi \lif \psi) \lor (\psi \lif \varphi)$
    \item $(\lnot\lnot \varphi \lif \varphi) \lif (\varphi \lor \lnot \varphi)$
    \item $(\varphi \lif \psi \lor \chi) \lif \bigl((\varphi \lif \psi)\lor(\varphi \lif \chi)\bigr)$
  \end{enumerate}
\end{prob}






\begin{quote}\small\textbf{स्रोतनोंद.} SOL6-A200: शून्य ते एक ही विगमन-पायरी स्पष्ट केली. सांत proof support रिकामा असल्यास शून्य टप्पा घेतला; नवीन सूत्र न मिळवता स्थिर राहणारी शाखा जपली. आधीचे fairness argument बदलले नाही.\end{quote}






\section{लिंडेनबाउमचे पूर्वप्रमेय}\label{int:sc:lin:sec}

अंतःप्रज्ञावादी तर्कशास्त्राचे संपूर्णता प्रमेय सिद्ध करताना
$\Gamma \Proves/ \varphi$ असे गृहीत धरून
$\mSat{M}{\Gamma}$ आणि~$\mSat/{M}{\varphi}$
असे प्रतिमान रचले जाते.

अभिजात तर्कशास्त्रात निष्पन्नता चा संबंध
सुसंगततेच्या संकल्पनेत आणता येतो. कारण सूत्रांच्या
एखाद्या संचापासून सूत्र~$\varphi$
निष्पन्न करता येण्याजोगे असते तेव्हा आणि केवळ तेव्हाच
त्या संचात~$\varphi$ चे नकरण घातल्यावर मिळणारा
संच विसंगत असतो. अंतःप्रज्ञावादी तर्कशास्त्रात
हे शक्य नाही. अंतःप्रज्ञावादी तर्कशास्त्रात
$\lnot\varphi$ विसंगत असल्यास केवळ
$\Proves \lnot\lnot \varphi$ हे मिळते.
$\lnot\lnot\varphi \lif \varphi$ हे सर्वसाधारणपणे
अंतःप्रज्ञावादी रीतीने खरे नसल्याने
$\Proves \varphi$ असा निष्कर्ष काढता येत नाही.

म्हणून प्रतिमान~$\mModel{M}$ रचताना
सूत्र~$\varphi$ ची निष्पन्नता नसणे
याची नोंद ठेवावी लागेल. त्यामुळे अभिजात
पद्धतीप्रमाणे प्रतिमान~$\mModel{M}$
रचण्यासाठी $\Gamma^* \supseteq \Gamma$
असा संपूर्ण संच वापरता येणार नाही;
कारण प्रत्येक संपूर्ण संच~$\Gamma^*$ साठी
$\Gamma^* \Proves \varphi \lor \lnot \varphi$.

संपूर्ण संच~$\Gamma^*$ वापरण्याऐवजी
सूत्रांच्या अभाज्य संचाची संकल्पना वापरू:

\begin{defn}\label{int:sc:lin:defn:prime}
  सूत्रांचा संच~$\Gamma$
  \emph{अभाज्य} असतो तेव्हा आणि केवळ तेव्हाच
  पुढील अटी पूर्ण होतात:
  \begin{enumerate}
  \item\label{int:sc:lin:defn:prime1} $\Gamma$ सुसंगत असतो,
    म्हणजे $\Gamma \Proves/ \lfalse$;
  \item\label{int:sc:lin:defn:prime2} $\Gamma \Proves \varphi$
    असेल, तर $\varphi \in \Gamma$; आणि
  \item\label{int:sc:lin:defn:prime3} $\varphi \lor \psi \in \Gamma$
    असेल, तर $\varphi \in \Gamma$ किंवा
    $\psi \in \Gamma$.
  \end{enumerate}
\end{defn}

\begin{lem}[लिंडेनबाउमचे पूर्वप्रमेय]
  \label{int:sc:lin:lem:lindenbaum}
  $\Gamma \Proves/ \varphi$ असेल, तर
  $\Gamma^* \supseteq \Gamma$ असा
  अभाज्य~$\Gamma^*$ अस्तित्वात असतो
  आणि $\Gamma^* \Proves/ \varphi$.
\end{lem}

\begin{proof}
  $\psi \lor \chi$ या रूपातील सर्व
  सूत्रांचे एक प्रगणन
  $\psi_1 \lor \chi_1$, $\psi_2 \lor \chi_2$, \dots,
  असू द्या. सूत्रांच्या संचांची
  वाढती श्रेणी~$\Gamma_n$ परिभाषित करू;
  प्रत्येक $\Gamma_{n+1}$ हा $\Gamma_n$
  मध्ये जास्तीतजास्त एक नवे सूत्र मिळवून
  परिभाषित केला जाईल.
  $\Gamma^*$ हा सर्व~$\Gamma_n$ चा
  संयोग असेल. नवी सूत्रे अशी निवडू
  की $\Gamma^*$ अभाज्य राहील आणि
  $\Gamma^* \Proves/ \varphi$ सुद्धा राहील.
  यासाठी प्रत्येक पायरीवर पुढील अटी
  पूर्ण करणारा पहिला विकल्पयोग~
  $\psi_i \lor \chi_i$ शोधला पाहिजे:
  \begin{enumerate}
  \item\label{int:sc:lin:gamma-1}
    $\Gamma_n \Proves \psi_i \lor \chi_i$
  \item\label{int:sc:lin:gamma-2}
    $\psi_i \notin \Gamma_n$ आणि
    $\chi_i \notin \Gamma_n$
  \end{enumerate}
  $\Gamma_n \cup \{\psi_i\} \Proves/ \varphi$
  असेल, तर $\Gamma_n$ मध्ये $\psi_i$
  मिळवू; अन्यथा $\chi_i$ मिळवू.
  हे यशस्वी होते ते दाखवायचे आहे.
  सध्या \ref{int:sc:lin:gamma-1} आणि~
  \ref{int:sc:lin:gamma-2} पूर्ण करणारा सर्वात
  लहान $i$ म्हणजे $i(n)$ असे ठरवू.

  $\Gamma_0 = \Gamma$ असे ठेवून
  \[  \Gamma_{n+1} = \begin{cases}
    \Gamma_n \cup \{\psi_{i(n)}\} &
    \text{जर $\Gamma_n \cup \{\psi_{i(n)}\} \Proves/ \varphi$ असेल} \\
    \Gamma_n \cup \{\chi_{i(n)}\} & \text{अन्यथा}
    \end{cases}
  \]
  अशी व्याख्या करू. $i(n)$ अपरिभाषित
  असेल, म्हणजे $\Gamma_n \Proves \psi \lor \chi$
  असताना नेहमी $\psi \in \Gamma_n$
  किंवा $\chi \in \Gamma_n$ असेल, तर
  $\Gamma_{n+1} = \Gamma_n$ ठेवू.
  आता $\Gamma^* = \bigcup_{n=0}^\infty \Gamma_n$
  असे ठेवू.

  प्रथम प्रत्येक $n$ साठी
  $\Gamma_n \Proves/ \varphi$ हे दाखवू.
  $n$ वर विगमन करू. $n = 0$ साठी
  हा दावा प्रमेयाच्या गृहीतकामुळे,
  म्हणजे $\Gamma \Proves/ \varphi$ मुळे,
  खरा आहे. कोणत्याही $n\geq0$ साठी
  $\Gamma_n \Proves/ \varphi$ असल्यास
  $\Gamma_{n+1} \Proves/ \varphi$
  हे दाखवायचे आहे.
  $i(n)$ अपरिभाषित असेल, तर
  $\Gamma_{n+1} = \Gamma_n$;
  काहीच सिद्ध करावे लागत नाही.
  म्हणून $i(n)$ परिभाषित आहे असे
  समजा. सोयीसाठी $i = i(n)$ असे लिहू.

  दाव्याचे व्यत्यस्त रूप सिद्ध करू.
  $\Gamma_{n+1} \Proves \varphi$ असे समजा.
  रचनेनुसार $\Gamma_n \cup
  \{\psi_i\} \Proves/ \varphi$ असेल, तर
  $\Gamma_{n+1} = \Gamma_n \cup
  \{\psi_i\}$; अन्यथा
  $\Gamma_{n+1} = \Gamma_n \cup \{\chi_i\}$.
  पहिला प्रकार असूच शकत नाही;
  कारण तेव्हा $\Gamma_{n+1} \Proves/ \varphi$.
  म्हणून $\Gamma_n \cup
  \{\psi_i\} \Proves \varphi$ आणि
  $\Gamma_{n+1} = \Gamma_n \cup \{\chi_i\}$.
  $i(n)$ च्या व्याख्येनुसार
  $\Gamma _n \Proves \psi_i \lor \chi_i$.
  $\Gamma_n \cup \{\psi_i\} \Proves \varphi$
  आहे. तसेच
  $\Gamma_{n+1} = \Gamma_n \cup
  \{\chi_i\} \Proves \varphi$ आहे.
  त्यामुळे $\Gamma_n \Proves \varphi$,
  हेच दाखवायचे होते.

  $\Gamma^* \Proves \varphi$ असेल, तर
  $\Gamma' \subseteq \Gamma^*$
  असा काही परिमित उपसंच असेल की
  $\Gamma' \Proves \varphi$.
  प्रत्येक $\theta \in \Gamma'$ कोणत्या
  तरी~$i$ साठी $\Gamma_i$ मध्ये असले
  पाहिजे. $\Gamma'$ रिकामा असेल, तर $n=0$ घ्या;
  अन्यथा या निर्देशांकांपैकी सर्वात मोठा $n$ घ्या. $i \le n$ असेल, तर
  $\Gamma_i \subseteq \Gamma_n$;
  म्हणून $\Gamma' \subseteq \Gamma_n$.
  पण मग $\Gamma_n \Proves \varphi$,
  आणि वर सिद्ध केलेल्या
  $\Gamma_n \Proves/ \varphi$ शी हा विरोधाभास.

  शेवटी $\Gamma^*$ अभाज्य आहे,
  म्हणजे \ref{int:sc:lin:defn:prime} मधील
  \ref{int:sc:lin:defn:prime1},
  \ref{int:sc:lin:defn:prime2} आणि~
  \ref{int:sc:lin:defn:prime3} या अटी
  पूर्ण करतो, हे दाखवू.

  प्रथम $\Gamma^* \Proves/ \varphi$;
  म्हणून $\Gamma^*$ सुसंगत आहे,
  आणि \ref{int:sc:lin:defn:prime1} पूर्ण होते.

  आता $\Gamma^* \Proves \psi \lor \chi$
  असेल, तर $\psi \in \Gamma^*$
  किंवा $\chi \in \Gamma^*$,
  हे दाखवू. यावरून
  \ref{int:sc:lin:defn:prime3} सिद्ध होते;
  कारण $\psi \lor \chi \in \Gamma^*$
  असेल, तर $\Gamma^* \Proves \psi
  \lor \chi$ सुद्धा. म्हणून
  $\Gamma^* \Proves \psi \lor \chi$
  पण $\psi \notin \Gamma^*$
  आणि $\chi \notin \Gamma^*$
  असे समजा.
  $\Gamma^* \Proves \psi \lor \chi$
  असल्याने कोणत्या तरी~$n$ साठी
  $\Gamma_n \Proves \psi \lor \chi$.
  $\psi \lor \chi$ हा विकल्पयोग
  प्रगणनात, समजा $\psi_j \lor \chi_j$
  म्हणून, येतो.
  $\psi_j \lor \chi_j$ हा
  $i(n)$ च्या व्याख्येतील गुणधर्म
  पूर्ण करतो: $\Gamma_n \Proves \psi_j
  \lor \chi_j$, तर
  $\psi_j \notin \Gamma_n$ आणि
  $\chi_j \notin \Gamma_n$.
  त्या निर्देशांकापेक्षा आधीचे निर्देशांक
  परिमित संख्येने आहेत. हा विकल्पयोग
  निवडला जात नाही तोवर अटी पूर्ण करत
  राहतो; तोवर निवडलेला विकल्पयोग
  $\psi_i \lor \chi_i$ त्याच्या आधीचाच असतो.
  त्यातील $\psi_i$ किंवा $\chi_i$
  मिळवल्यावर तो पुन्हा अटी पूर्ण करत नाही.
  म्हणून काही टप्प्यावर~$m$
  असे मिळेल की $j = i(m)$.
  पण तेव्हा $\psi \in \Gamma_{m+1}$
  किंवा $\chi \in \Gamma_{m+1}$;
  आणि हे $\psi \notin \Gamma^*$
  व $\chi \notin \Gamma^*$ या
  गृहीतकाच्या विरुद्ध आहे.

  आता $\Gamma^* \Proves \psi$
  असे समजा. मग
  $\Gamma^* \Proves \psi \lor \psi$.
  पण $\Gamma^* \Proves \psi \lor \psi$
  असेल, तर $\psi \in \Gamma^*$,
  हे आत्ताच सिद्ध केले आहे.
  त्यामुळे $\Gamma^*$ हा
  \ref{int:sc:lin:defn:prime} मधील
  \ref{int:sc:lin:defn:prime2} पूर्ण करतो.
\end{proof}

\begin{prob}
$\Gamma \Proves/ \bot$ असेल, तर
अभिजात तर्कशास्त्रात $\Gamma$ सुसंगत
असतो, म्हणजे~$\Gamma$ मधील सर्व सूत्रांना
सत्य करणारे मूल्यांकन अस्तित्वात असते,
हे दाखवा.
\end{prob}






\begin{quote}\small\textbf{स्रोतनोंद.} SOL6-A201: मूळ one-based pair list आणि शून्यासहित natural-number worlds यांतील विसंगती दुरुस्त केली. सर्व नैसर्गिक निर्देशांकांसाठी formula-pair enumeration शून्यापासून सुरू केली; शाखा आणि truth-lemma साठी लागणाऱ्या गुणधर्मांना बदल नाही.\end{quote}






\section{कॅनॉनिकल प्रतिमान}\label{int:sc:mod:sec}

आपल्या प्रतिमानातील जगे नैसर्गिक संख्यांच्या
सांत क्रमिका~$\sigma$ असतील, म्हणजे
$\sigma \in \Nat^*$. $\Nat^*$ ची
विगमनाने पुढीलप्रमाणे व्याख्या होते:
\begin{enumerate}
\item $\emptyseq \in \Nat^*$.
\item $\sigma \in \Nat^*$ आणि $n \in \Nat$
  असतील, तर $\sigma.n \in
  \Nat^*$ (येथे $\sigma.n$ म्हणजे
  $\sigma \concat \tuple{n}$, आणि
  $\sigma \concat \sigma'$ म्हणजे
  $\sigma$ व $\sigma'$ यांची जोडणी).
\item इतर काहीही $ \Nat^*$ मध्ये नाही.
\end{enumerate}
म्हणून $\Nat^*$ वापरून विगमनात्मक
व्याख्या देता येतात.

सूत्रांच्या सर्व जोड्यांचे
$\tuple{\psi_0, \chi_0}$,
$\tuple{\psi_1, \chi_1}$,
$\tuple{\psi_2, \chi_2}$, \dots,
असे $n\in\Nat=\{0,1,2,\dots\}$ ने निर्देशांकित
प्रगणन असू द्या. सूत्रे
चा संच~$\Delta$ दिला असता
$\Delta(\sigma)$ ची विगमनाने
पुढीलप्रमाणे व्याख्या करा:
\begin{enumerate}
\item $\Delta(\emptyseq) = \Delta$
\item $\Delta(\sigma.n) = {}$
  \[  \begin{cases}
    (\Delta(\sigma) \cup \{\psi_n\})^* &
    \text{जर $\Delta(\sigma) \cup \{\psi_n\} \Proves/ \chi_n$ असेल} \\
    \Delta(\sigma) & \text{अन्यथा}
  \end{cases}
  \]
\end{enumerate}
येथे $(\Delta(\sigma) \cup \{\psi_n\})^*$
म्हणजे $\Delta(\sigma) \cup \{\psi_n\}$
या संचाला आणि सूत्र~$\chi_n$
ला \ref{int:sc:lin:lem:lindenbaum} लावल्यावर
मिळणारा सूत्रांचा अभाज्य संच.
या व्याख्येनुसार
$\Delta(\sigma) \cup \{\psi_n\} \Proves/ \chi_n$
असेल, तर $\Delta(\sigma.n)
\Proves \psi_n$ आणि
$\Delta(\sigma.n) \Proves/ \chi_n$,
हे लक्षात घ्या. तसेच कोणत्याही~$n$
साठी $\Delta(\sigma) \subseteq \Delta(\sigma.n)$.
$\Delta$ अभाज्य असेल, तर सर्व~$\sigma$
साठी $\Delta(\sigma)$ अभाज्य असतो.

\begin{defn}\label{int:sc:mod:defn:canonical-model}
  $\Delta$ अभाज्य आहे असे समजा.
  मग $\Delta$ साठीचे
  \emph{कॅनॉनिकल प्रतिमान}
  $\mModel{M(\Delta)}$ पुढीलप्रमाणे परिभाषित होते:
  \begin{enumerate}
  \item $W = \Nat^*$, म्हणजे नैसर्गिक
    संख्यांच्या सांत क्रमिकांचा संच.
  \item $R$ हा असा अंशतः क्रम आहे की
    $R\sigma\sigma'$ तेव्हा आणि केवळ
    तेव्हाच जेव्हा $\sigma$ हा~$\sigma'$
    चा आरंभीचा खंड असतो (म्हणजे
    काही क्रमिका~$\sigma''$ साठी
    $\sigma' = \sigma \concat \sigma''$).
  \item $V(p) = \Setabs{\sigma}{p \in \Delta(\sigma)}$.
  \end{enumerate}
\end{defn}

$R$ हा खरोखर अंशतः क्रम आहे,
हे सहज तपासता येते. तसेच~$V$ वरील
एकस्वनिकतेची अट पूर्ण होते.
$\Delta(\sigma) \subseteq \Delta(\sigma.n)$
असल्याने $R\sigma\sigma'$ असेल तेव्हा
$\Delta(\sigma) \subseteq \Delta(\sigma')$;
हे~$\sigma$ पासूनच्या विस्ताराच्या
लांबीवर विगमनाने मिळते.












\section{सत्यता पूर्वप्रमेय}\label{int:sc:tru:sec}

पुढील पूर्वप्रमेय कॅनॉनिकल प्रतिमानातील
पूर्तीचा संबंध अभाज्य संच
$\Delta(\sigma)$ मध्ये कोणती
सूत्रे घटक आहेत
याच्याशी जोडते.

\begin{lem}\label{int:sc:tru:lem:truth}
  $\Delta$ अभाज्य असेल, तर
  $\mSat{M(\Delta)}{\varphi}[\sigma]$
  तेव्हा आणि केवळ तेव्हाच
  $\Delta(\sigma) \Proves \varphi$.
\end{lem}

\begin{proof}
  $\varphi$ वर विगमन करू.
  \begin{enumerate}
    \item \indcase{\varphi}{\lfalse}{$\Delta(\sigma)$
      अभाज्य असल्याने सुसंगत आहे;
      म्हणून $\Delta(\sigma) \Proves/ \indfrm$.
      व्याख्येनुसार
      $\mSat/{M(\Delta)}{\indfrm}[\sigma]$.}
    \item \indcase{\varphi}{p}{$\mSat{{}}{}$ च्या
        व्याख्येनुसार
        $\mSat{M(\Delta)}{\indfrm}[\sigma]$
        तेव्हा आणि केवळ तेव्हाच
        $\sigma \in V(p)$;
        म्हणजे $\Delta(\sigma) \Proves \indfrm$;
        कारण हा अभाज्य संच निष्पत्तीखाली संवृत आहे.}
    \item \indcase!{\varphi}{\lnot \psi}{}
    \item \indcase{\varphi}{\psi \land
        \chi}{$\mSat{M(\Delta)}{\indfrm}[\sigma]$
        तेव्हा आणि केवळ तेव्हाच
        $\mSat{M(\Delta)}{\psi}[\sigma]$
        आणि $\mSat{M(\Delta)}{\chi}[\sigma]$.
        विगमन-गृहीतकानुसार
        $\mSat{M(\Delta)}{\psi}[\sigma]$
        तेव्हा आणि केवळ तेव्हाच
        $\Delta(\sigma) \Proves
        \psi$; आणि~$\chi$ साठीही तसेच.
        पण $\Delta(\sigma) \Proves \psi$
        आणि $\Delta(\sigma) \Proves \chi$
        हे तेव्हा आणि केवळ तेव्हाच,
        जेव्हा $\Delta(\sigma)
        \Proves \indfrm$.}
    \item \indcase{\varphi}{\psi \lor
        \chi}{$\mSat{M(\Delta)}{\indfrm}[\sigma]$
        तेव्हा आणि केवळ तेव्हाच
        $\mSat{M(\Delta)}{\psi}[\sigma]$
        किंवा $\mSat{M(\Delta)}{\chi}[\sigma]$.
        विगमन-गृहीतकानुसार हे
        तेव्हा आणि केवळ तेव्हाच,
        जेव्हा $\Delta(\sigma) \Proves \psi$
        किंवा $\Delta(\sigma) \Proves \chi$.
        आता हे पुन्हा
        $\Delta(\sigma) \Proves \indfrm$
        च्या तुल्य आहे, हे दाखवायचे आहे.
        डावीकडून उजवीकडील दिशा उघड आहे.
        उजवीकडून डावीकडील दिशा
        $\Delta(\sigma)$ अभाज्य
        असल्याने मिळते.}
    \item \indcase{\varphi}{\psi \lif \chi}{आधी
        डावीकडून उजवीकडील दिशेचे
        व्यत्यस्त रूप पाहू:
        $\Delta(\sigma) \Proves/ \psi
        \lif \chi$ असे समजा.
        मग $\Delta(\sigma) \cup \{\psi\} \Proves/
        \chi$ सुद्धा.
        कोणत्या तरी~$n$ साठी
        $\tuple{\psi, \chi}$ ही
        $\tuple{\psi_n, \chi_n}$ असल्याने
        $\Delta(\sigma.n) = (\Delta(\sigma) \cup
        \{\psi\})^*$; आणि
        $\Delta(\sigma.n) \Proves \psi$
        पण $\Delta(\sigma.n) \Proves/
        \chi$. विगमन-गृहीतकानुसार
        $\mSat{M(\Delta)}{\psi}[\sigma.n]$
        आणि $\mSat/{M(\Delta)}{\chi}[\sigma.n]$.
        $R\sigma(\sigma.n)$ असल्यामुळे
        याचा अर्थ
        $\mSat/{M(\Delta)}{\indfrm}[\sigma]$.

        आता $\Delta(\sigma) \Proves \psi \lif \chi$
        असे समजा आणि $R\sigma\sigma'$ असू द्या.
        $\Delta(\sigma) \subseteq
        \Delta(\sigma')$ असल्याने
        $\Delta(\sigma') \Proves
        \psi$ असेल, तर
        $\Delta(\sigma') \Proves \chi$.
        दुसऱ्या शब्दांत, $R\sigma\sigma'$
        असलेल्या प्रत्येक $\sigma'$ साठी
        $\Delta(\sigma') \Proves/ \psi$
        किंवा $\Delta(\sigma') \Proves
        \chi$. विगमन-गृहीतकानुसार
        $R\sigma\sigma'$ असेल तेव्हा
        $\mSat/{M(\Delta)}{\psi}[\sigma']$
        किंवा $\mSat{M(\Delta)}{\chi}[\sigma']$;
        म्हणजे
        $\mSat{M(\Delta)}{\indfrm}[\sigma]$.}
  \end{enumerate}
\end{proof}












\section{संपूर्णता प्रमेय}\label{int:sc:cpl:sec}

\begin{thm}\label{int:sc:cpl:thm:completeness}
  $\Gamma \Entails \varphi$ असेल, तर
  $\Gamma \Proves \varphi$.
\end{thm}

\begin{proof}
  व्यत्यस्त रूप सिद्ध करू:
  $\Gamma \Proves/ \varphi$ असे समजा.
  मग \ref{int:sc:lin:lem:lindenbaum} नुसार
  $\Gamma^* \supseteq \Gamma$ असा
  अभाज्य संच आहे की
  $\Gamma^* \Proves/ \varphi$.
  \ref{int:sc:mod:defn:canonical-model}
  मध्ये व्याख्यित केलेले $\Gamma^*$ साठीचे
  कॅनॉनिकल प्रतिमान~$\mModel{M(\Gamma^*)}$
  घ्या. प्रत्येक $\psi \in \Gamma$
  साठी $\Gamma^* \Proves \psi$.
  $\Gamma^*(\emptyseq) = \Gamma^*$,
  हे लक्षात घ्या. सत्यता पूर्वप्रमेयानुसार
  (\ref{int:sc:tru:lem:truth})
  सर्व $\psi \in \Gamma$ साठी
  $\mSat{M(\Gamma^*)}{\psi}[\emptyseq]$
  आणि
  $\mSat/{M(\Gamma^*)}{\varphi}[\emptyseq]$.
  यावरून $\Gamma \Entails/ \varphi$.
\end{proof}

\begin{prob}
  $\varphi$ मध्ये फक्त विधानीय चले,
  $\lor$ आणि~$\land$ असतील, तर
  $\Entails/ \varphi$ हे दाखवा.
  यावरून अंतःप्रज्ञावादी तर्कशास्त्रात
  $\lor$ आणि~$\land$ वापरून
  $\lif$ ची व्याख्या करता येत नाही,
  असा निष्कर्ष काढा.
\end{prob}

\begin{prob}
  संपूर्णता प्रमेय वापरून
  $\Proves \varphi \lor \psi$ असेल,
  तर $\Proves \varphi$ किंवा
  $\Proves \psi$, हे सिद्ध करा.
  (सूचना: $\mSat/{M_1}{\varphi}$
  आणि $\mSat/{M_2}{\psi}$
  असे गृहीत धरून
  $\mSat/{M}{\varphi \lor \psi}$
  असेल असे नवे प्रतिमान
  $\mModel{M}$ रचा.)
\end{prob}

\begin{prob}
  $\mModel{M}$ हे रेषीय क्रम
  वापरणारे संबंधी प्रतिमान असेल,
  तर $\mSat{M}{(\varphi \lif \psi)\lor(\psi \lif \varphi)}$,
  हे दाखवा.
\end{prob}






\begin{quote}\small\textbf{स्रोतनोंद.} SOL6-A203: मूळ atom-only quotient चुकीचा असल्याची OLINC-226 सूचना योग्य होती; तरी त्याच अवैध truth-equivalence ची सिद्धता exercise मध्ये मागितली होती. आता सर्व उपसूत्रांवरील योग्य गाळणी, implication आणि primitive negation सहित truth-preservation proof, finite bound आणि bounded decision procedure दिली. Exercise योग्य गाळणीच्या तपशिलांसाठी केला.\end{quote}






\section{निर्णेयता}\label{int:sc:dec:sec}
संपूर्णता प्रमेयाची सिद्धता
$\Gamma \Proves/ \varphi$ असलेल्या
प्रत्येक वेळी $\Gamma \Entails/ \varphi$
याचा साक्षीदार म्हणून अनंत जगे
असलेले प्रतिमान देते, हे लक्षात घ्या.
पुढील विधानानुसार
$\Entails \varphi$ सिद्ध करण्यासाठी
सर्व सांत प्रतिमानांमध्ये (म्हणजे
जगांचा सांत संच असलेल्या प्रतिमानांमध्ये)
$\mSat{M}{\varphi}$ हे सिद्ध करणे पुरेसे आहे.

\begin{thm}\label{int:sc:dec:thm:decidability}
  $\Entails/ \varphi$ असेल, तर
  $\mSat/{M'}{\varphi}$ असे सांत प्रतिमान
  अस्तित्वात असते.
\end{thm}

\begin{proof}
  $\mModel{M}=\tuple{W,R,V}$ हे $\mSat/{M}{\varphi}$
  असे प्रतिमान असू द्या. $S$ हा $\varphi$ च्या सर्व उपसूत्रांचा
  सांत संच असू द्या; विशेषतः $\varphi\in S$ आणि $S$
  उपसूत्रे घेण्याखाली संवृत आहे. $P$ हा $\varphi$ मध्ये
  येणाऱ्या विधानीय चलांचा संच आहे.
  प्रत्येक $w\in W$ साठी
  \[    [w]=\Setabs{\psi\in S}{\mSat{M}{\psi}[w]}
  \]
  असे ठरवून $\mModel{M'}=\tuple{W',R',V'}$ परिभाषित करू:
  \[    W'=\Setabs{[w]}{w\in W},\qquad
    R'[w][u]\ \text{तेव्हा आणि केवळ तेव्हाच}\ [w]\subseteq[u],
  \]
  आणि $p\in P$ साठी $V'(p)=\Setabs{[w]\in W'}{p\in[w]}$;
  इतर विधानचरांसाठी $V'(p)=\emptyset$.
  $W'$ अरिक्त असून $|W'|\leq 2^{|S|}$.
  उपसंच-संबंध हा अंशतः क्रम आहे, आणि $V'$ ची
  एकस्वनिकता व्याख्येवरून मिळते. म्हणून $\mModel{M'}$
  हे सांत संबंधी प्रतिमान आहे.

\begin{quote}\small\textbf{स्रोतनोंद.} OLINC-226: मूळ atom-only quotient उपसूत्रांची सत्यता जपत नव्हता. खाली त्याच्या जागी सर्व उपसूत्रांवरील योग्य गाळणीची सिद्धता आहे; समतुल्यता S मधील सूत्रांसाठीच सांगितली आहे.\end{quote}
  प्रत्येक $\psi\in S$ साठी रचनेवरील विगमनाने
  \[    \mSat{M}{\psi}[w]\ \text{तेव्हा आणि केवळ तेव्हाच}\
    \mSat{M'}{\psi}[{[w]}]
  \]
  हे सिद्ध करू. विधानचरांसाठी दावा $V'$ च्या व्याख्येवरून
  मिळतो; $\lfalse$ दोन्ही प्रतिमानांत कोठेही सत्य नसतो.
  $\land$ आणि $\lor$ यांचे प्रकार त्यांच्या स्थानिक
  सत्यता-अटी आणि विगमन-गृहीतकावरून मिळतात.
  $\psi\ident \chi\lif \theta$ असू द्या. मूळ प्रतिमानात
  $\mSat{M}{\psi}[w]$ असेल आणि $[w]\subseteq[u]$,
  तर $\psi\in[w]\subseteq[u]$; म्हणून $\mSat{M}{\psi}[u]$.
  $\mSat{M'}{\chi}[{[u]}]$ असेल, तर विगमन-गृहीतकामुळे
  $\mSat{M}{\chi}[u]$; मूळ संबंधाच्या परावर्तितेने
  $\mSat{M}{\theta}[u]$, आणि पुन्हा विगमन-गृहीतकामुळे
  $\mSat{M'}{\theta}[{[u]}]$. म्हणून $\mSat{M'}{\psi}[{[w]}]$.
  उलट, $\mSat/{M}{\psi}[w]$ असेल, तर काही $u$ साठी
  $Rwu$, $\mSat{M}{\chi}[u]$ आणि $\mSat/{M}{\theta}[u]$.
  मूळ प्रतिमानातील सत्यतेच्या एकस्वनिकतेमुळे
  $[w]\subseteq[u]$; विगमन-गृहीतकाने $[u]$ ही नव्या
  प्रतिमानातही त्या सशर्त सूत्राची प्रतिउदाहरण-अवस्था आहे.
  $\psi\ident\lnot \chi$ असेल, तर त्याच पद्धतीने:
  $\mSat{M}{\psi}[w]$ आणि $[w]\subseteq[u]$ वरून
  $\mSat{M}{\psi}[u]$, म्हणून $\mSat/{M}{\chi}[u]$ आणि
  $\mSat/{M'}{\chi}[{[u]}]$. त्यामुळे नव्या प्रतिमानातही
  $\lnot \chi$ सत्य आहे. मूळ प्रतिमानात $\lnot \chi$ असत्य
  असेल, तर $Rwu$ आणि $\mSat{M}{\chi}[u]$ असलेली अवस्था
  $u$ मिळते; तिच्यासाठी $[w]\subseteq[u]$ आणि
  $\mSat{M'}{\chi}[{[u]}]$, म्हणून नव्या प्रतिमानातही नकरण
  असत्य आहे. प्रत्येक प्रकारात आवश्यक घटक-उपसूत्रे $S$
  मध्ये असल्याने विगमन-गृहीतक लागू आहे.
  अखेरीस $\mSat/{M}{\varphi}$ असल्याने काही $w\in W$
  साठी $\mSat/{M}{\varphi}[w]$; समतुल्यतेवरून
  $\mSat/{M'}{\varphi}[{[w]}]$. म्हणून $\mSat/{M'}{\varphi}$.
\end{proof}

\begin{prob}
\ref{int:sc:dec:thm:decidability} मधील योग्य
उपसूत्र-गाळणीचे तपशील तपासा: $V'$ नेमके परिभाषित आहे,
अंशतः क्रमाच्या दृष्टीने एकस्वनिक आहे, आणि $|W'|\leq2^{|S|}$
हे दाखवा. $S$ उपसूत्रे घेण्याखाली संवृत असण्याची अट
विगमन-सिद्धतेत कुठे वापरली आहे ते स्पष्ट करा.
\end{prob}

\ref{int:sc:dec:thm:decidability} आणि त्याच्या सिद्धतेतील
मर्यादेवरून $\Entails \varphi$ ठरवणारा अल्गोरिदम मिळतो.
$S$ आणि $P$ प्रभावीपणे मोजा. $1\leq n\leq2^{|S|}$
असलेल्या प्रत्येक $n$ साठी $n$ जगांवरील सर्व अंशतः क्रम
आणि $P$ वरील सर्व एकस्वनिक मूल्यांकने यांचा सांत शोध करा.
इतर विधानचरांना रिकामा संच दिल्याने $\varphi$ ची सत्यता बदलत नाही.
प्रत्येक अशा प्रतिमानातील प्रत्येक जगात $\varphi$ च्या सत्यतेचे
recursive clauses वापरून मूल्यांकन करा. प्रतिउदाहरण मिळाल्यास
$\Entails/ \varphi$; कोणतेही मिळाले नाही, तर प्रमेयानुसार
$\Entails \varphi$. हा शोध नेहमी संपतो.




